\section{Preliminaries and Problem Formulation}
\label{sec:definition}

\paragraph{Coding-agent episodes.}
Let $q$ denote a repository task and let an agent policy $\pi$ interact with a workspace through tool calls, shell commands, file edits, and tests. A one-shot agent produces a trajectory
\begin{equation}
\tau \sim \pi(\cdot \mid q),
\end{equation}
where $\tau$ contains observations, actions, tool outputs, and edits. For short tasks, a single trajectory may be enough. For long-horizon engineering, however, the desired object is not one trajectory but a sequence of recoverable epochs.

\begin{definition}[Reflective loop state]
A reflective loop state is a tuple
\begin{equation}
\mathcal{L}^{(k)} = \langle R, A^{(k)}, F^{(k)}, M^{(k)}, H^{(k)} \rangle,
\end{equation}
where $R$ is the roadmap, $A^{(k)}$ is the active task adapter, $F^{(k)}$ is the failure memory, $M^{(k)}$ records the current mode, and $H^{(k)}$ stores logs, snapshots, and handoff notes at epoch $k$.
\end{definition}

\paragraph{Slow and fast state.}
The roadmap $R$ stores durable project intent: goals, constraints, milestones, and non-negotiable design choices. It is analogous to a frozen or slowly updated backbone. The active adapter $A^{(k)}$ is a structured JSON record for the current slice of work. It stores local patches, immediate verification requirements, and the next bounded target. This separation prevents two common failures: global plans being overwritten by local errors, and local fixes being lost inside a long transcript.

\paragraph{Optimize/check epochs.}
Each epoch contains two conceptual phases. The optimize phase performs bounded implementation conditioned on $(R,A^{(k)},F^{(k)})$. The check phase reviews the result, estimates the residual error, and emits patch signals:
\begin{equation}
P^{(k)} = C(\tau^{(k)}, R, A^{(k)}, F^{(k)}),
\end{equation}
where $P^{(k)}$ may include prompt patches, scope patches, verification patches, and new failure patterns. The next state is updated externally:
\begin{equation}
A^{(k+1)}, F^{(k+1)}, H^{(k+1)} = U(A^{(k)}, F^{(k)}, H^{(k)}, P^{(k)}).
\end{equation}
This update is non-parametric: it changes files and structured state, not model weights.

\paragraph{Objective.}
Let $S(\tau,q)$ be a task success indicator and let $B(\tau)$ measure engineering cost, such as repeated failed attempts or redundant context reconstruction. \ours{} aims to improve long-horizon execution by increasing expected success while reducing repeated work:
\begin{equation}
\max_{\mathcal{L}} \; \mathbb{E}_{q}\left[S(\tau,q) - \lambda B(\tau)\right],
\end{equation}
subject to auditability: every epoch should leave enough state for a later agent or human to resume the project without relying on hidden context.

